\section{Project memory, roadmaps and issue tracking}

\subsection{Why a substantial project needs external memory}

C-LARA-2 is now a substantial evolving repository: the July 2026 snapshot
contains approximately 82,740 tracked lines accumulated over about nine
months. No participant, human or AI, can reliably retain all relevant
architectural decisions, experiments, operational procedures, unresolved
problems, user requirements and publication claims in conversational working
memory. The project therefore treats the repository as more than a container
for code. It is a version-controlled external memory that records current
behaviour, past decisions, intended direction and the procedures through which
people can inspect or change the system.

This external memory must be searchable, structured, linked to the code it
describes and maintained as the project changes. It is useful to human
collaborators, who need to understand and redirect work without becoming
full-time programmers or system administrators. It is equally useful to
Codex, which can inspect the same records when implementing a change,
diagnosing a failure, planning an experiment or preparing a report. The aim is
not to replace human judgement with accumulated text, but to make project
history and current constraints available for human review and AI-supported
operational work.

\subsection{Roadmaps and how-to documentation}

Roadmap files have been part of C-LARA-2 from the start. They record why a
feature area matters, describe current architecture and design choices,
preserve the history of false starts or revisions, identify priorities and
dependencies, and link to relevant issues, tests, commands and operational
procedures. A roadmap is consequently not a static feature wish list. It is a
maintained account of how a part of the project reached its current state and
what work should happen next.

How-to files translate that memory into procedures that human collaborators
can follow. They cover development setup, backups, deployment, server
administration, migration and diagnostics. These documents are especially
important in an AI-centred project because a human participant may need to
perform an action in a production environment while relying on Codex to
inspect repository context, explain the reason for the action and formulate
safe, concrete steps. The documents preserve that operational knowledge for
the next occasion rather than leaving it only in a transient conversation.

\subsection{Issues as AI-maintained project records}

Focused issue records were added in early May 2026. They can be entered and
updated by human collaborators through a tab in the online platform, but they
are not simply conventional tickets in a separate tracking system. Codex also
reads, revises and cross-links their repository representations as part of
ordinary project maintenance. An issue can therefore connect a human
observation or suggestion to a roadmap, code change, test, operational note
and later status update.

This makes the records semantically richer than a minimal ``report, assign,
close'' workflow. An issue is a maintained account of a bounded problem and
its relationship to the project memory. It also creates a responsibility:
AI-maintained records can become stale, misunderstand a human concern or
overstate completion unless humans review their framing. The point is not that
conventional issue trackers are unnecessary, but that the issue record can
participate directly in the technical and documentary work needed to resolve
it.

\subsection{Operational documentation and AWS work}

The most important how-to material concerns use of the AWS server. None of the
human participants is a specialist system administrator, but Codex has been
able to guide them through nontrivial server tasks by turning repository and
deployment context into human-executable diagnostic and repair procedures.
Humans perform and supervise the actions, retain responsibility for production
access and decide whether a proposed change is acceptable; Codex supplies much
of the operational interpretation and step-by-step guidance.

One important example is the installation and diagnosis of the read-only
sandbox later used by the project-understanding Assistant. The technical
function of the Assistant is discussed in Section~9. The relevant point here
is methodological: Codex could inspect the deployment context, propose
concrete procedures, and preserve the resulting operational knowledge in
version-controlled documentation. This allowed a team without prior sysadmin
specialism to carry out work that would otherwise have been difficult to
plan, explain and repeat.

\subsection{Current scale and limits}

\begin{table}[htbp]
\centering
\caption{Selected project-memory artefacts in the July 2026 repository snapshot.}
\begin{tabular}{lr}
\toprule
Artefact type & Count \\
\midrule
Roadmap files & 31 \\
How-to files & 11 \\
Canonical issue records & 41 \\
Tracked repository lines & 82,740 \\
\bottomrule
\end{tabular}
\end{table}

These figures do not measure software quality, and more documentation does not
automatically produce coherence. The records must be reviewed, corrected and
kept aligned with the code and deployment state they describe. Their value is
that they make a larger amount of project history, operational knowledge and
planned work available for inspection and redirection. In that sense, the
roadmaps, how-to files and issues are not secondary administrative outputs;
they are part of the infrastructure through which C-LARA-2 remains intelligible
and revisable as it grows. They also provide an evidential basis for the
repository-grounded Assistant discussed later in the report.
